\section{Introduction}
\label{sec:introduction}

The source-vortex panel method is a numerical technique used in aerodynamics to approximate the potential flow around
bodies of arbitrary shape, including those that generate lift.
It belongs to the family of boundary element methods, where the surface of a body is discretized into a finite number
of panels.
Unlike the source panel method, which can only represent non-lifting flows, or the vortex panel method alone, the
source-vortex panel method assigns each panel both a source distribution and a vortex distribution.
The source strengths are typically permitted to vary from panel to panel in order to enforce the impermeability boundary
condition on the body surface, while a single vortex strength is often shared across all panels to represent the
overall circulation required to generate lift.
The method is based on the assumptions of incompressible, irrotational and inviscid flow, allowing the velocity field
to be represented by a scalar potential that satisfies Laplace's equation.
By superimposing the effects of the source panels, the vortex panels, and the free-stream flow, the velocity
potential around the body can be computed.
As with the vortex panel method, the circulation is not arbitrary but is instead constrained by the Kutta condition,
which enforces smooth flow leaving the trailing edge and ensures a unique, physically realistic solution.
Once the velocity distribution is known, aerodynamic quantities such as the pressure coefficient and lift can be
obtained using Bernoulli's equation and the Kutta-Joukowski theorem, respectively.
By combining the thickness representing capability of source panels with the lift generating capability of vortex
panels, the source-vortex panel method is capable of predicting both the pressure distribution and the lift of
practical lifting bodies, such as airfoils, and represents a natural synthesis of the two simpler panel formulations.
